\chapter{The transcendental non-Archimedean Calabi--Yau theorem}
\index{Calabi--Yau theorem!transcendental non-Archimedean}\index{Monge--Ampere equation!transcendental}
\label{chap:mpwn-cy}

The model theory of \Cref{chap:mp-psh} supplies Monge--Amp\`ere measures and energy,
but it does not by itself solve
\[
                          \MA_\alpha(\phi)=\mu.
\tag{61.1}\label{eq:mpwn-ma-equation}
\]
The variational solution requires compactness, continuous envelopes,
orthogonality, and control of both potential-side and measure-side strong
topologies.  These ingredients were established in the transcendental
analytic setting by Mesquita-Piccione and Witt Nystr\"om
\cite{MPWN25}.

Throughout, $X$ is a connected compact K\"ahler manifold, $\dim_{\mathbf C}X=n$,
$\alpha$ is a K\"ahler class, $A$ denotes its pullback system on all
analytic models, and $V=\int_X\alpha^n$.  We use probability
Monge--Amp\`ere measures.  Thus no hidden factor $V$ occurs in
\eqref{eq:mpwn-ma-equation}.

There are three nearby frameworks, with different ambient hypotheses.  The
projective trivially valued theory of \cite{BJ22} supplies the algebraic
variational architecture.  The synthetic theory of
\cite{BJ25FiniteEnergy} isolates the energy estimates that do not depend on
a particular model category.  The present chapter uses the analytic-model
theorems of \cite{MPWN25} to verify the missing geometric inputs for an
arbitrary compact K\"ahler class.  References to the first two frameworks
below identify antecedents and reusable estimates; they do not replace the
analytic envelope and restricted-volume arguments.

\section{Weak compactness and non-pluripolar points}
\label{sec:mpwn-weak-compactness}

The weak topology on $\PSH(\alpha)$ is pointwise convergence on
$X^{\divis}$.  Since psh functions are finite at divisorial points and
determined there, this topology is Hausdorff.  Put
\[
 \PSH_{\sup}(\alpha)
 :=\{\phi\in\PSH(\alpha)\mid\sup_{X^{\bethan}}\phi=0\}.
\tag{61.2}\label{eq:mpwn-sup-normalized-psh}
\]
The supremum equals $\phi(v_{\triv})$, a fact inherited from the model
maximum principle.

\begin{theorem}\label{thm:mpwn-weak-compactness}
The space $\PSH_{\sup}(\alpha)$ is compact for the weak topology.
\end{theorem}

\begin{proof}[Imported envelope core, with the net correction]
For every divisorial $v$, normalized psh functions are uniformly bounded
at $v$.  Embed $\PSH_{\sup}(\alpha)$ in the product of these compact
coordinate intervals and use Tychonoff's theorem.  Given a net
$(\phi_i)$, first pass to a subnet $(\phi_{i_j})$ that converges in the
ambient product.  Form, over tails of this convergent subnet, the upper
envelopes
\[
              \psi_j:=\left(\sup_{k\succeq j}\phi_{i_k}\right)^* .
\tag{61.3}\label{eq:mpwn-tail-envelope}
\]
The upper-semicontinuous regularized supremum of a uniformly
upper-bounded family of psh functions is psh, and it agrees with the
pointwise supremum on $X^{\divis}$.  The $\psi_j$ decrease to a
sup-normalized psh function having the prescribed divisorial values.
Hence the image is closed in the product and therefore compact.
The imported core is the arbitrary-family envelope statement
\cite[Proposition~3.1.2]{MPWN25}.  This proves
\cite[Theorem~3.1.1]{MPWN25}.  The source writes ``subsequence'' after
invoking Tychonoff; in the possibly nonmetrizable topology of the theorem,
the convergent subnet argument just given is required.
\end{proof}

\begin{lemma}\label{lem:mpwn-energy-weak-upper}
The energy $E_\alpha$ is upper semicontinuous on
$\PSH_{\sup}(\alpha)$ for the weak topology.  Consequently, for
each $R\geq0$, the set
\[
 \{\phi\in\PSH_{\sup}(\alpha)\mid E_\alpha(\phi)\geq-R\}
\]
is weakly compact and is contained in $\mathcal E^1(\alpha)$.
\end{lemma}

\begin{proof}
Let $\phi_i$ converge weakly to $\phi$ in $\PSH_{\sup}(\alpha)$
and form the tail envelopes $\psi_i$ as in
\eqref{eq:mpwn-tail-envelope}.  They decrease to $\phi$:
their divisorial values converge to those of $\phi$, and
decreasing stability and divisorial determination identify the
limit.  For every $j\succeq i$, monotonicity gives
$E_\alpha(\psi_i)\geq E_\alpha(\phi_j)$.  Decreasing continuity
of the extended energy, including the value $-\infty$, yields
\[
 E_\alpha(\phi)=\inf_i E_\alpha(\psi_i)
                 \geq\limsup_i E_\alpha(\phi_i).
\]
The asserted set is therefore closed in the compact space of
\Cref{thm:mpwn-weak-compactness}.  Its energy lower bound gives
membership in $\mathcal E^1(\alpha)$.
\end{proof}


\begin{definition}[Non-pluripolar points]
\label{def:mpwn-nonpluripolar-point}
A point $v\in X^{\bethan}$ is non-pluripolar if
\[
             \phi(v)>-\infty\quad\text{for every }\phi\in\PSH(\alpha).
\tag{61.4}\label{eq:mpwn-nonpluripolar}
\]
Write $X^{\mathrm{np}}$ for this locus.  On it define
\[
 d_\infty(v,w)
   :=\sup_{\phi\in\PSH(\alpha)}
       |\phi(v)-\phi(w)|,
\qquad
 T(v):=d_\infty(v,v_{\triv}).
\tag{61.5}\label{eq:mpwn-dinfty}
\]
\end{definition}

Every divisorial point is non-pluripolar by
\Cref{thm:mp-psh-divisorial}.
\par\noindent\begin{minipage}{\linewidth}
For $v\in X^{\mathrm{np}}$, the maximum principle gives
\[
            T(v)=\sup_{\phi\in\PSH_{\sup}(\alpha)}(-\phi(v)).
\]
\end{minipage}\par
To prove finiteness, suppose this supremum is infinite and choose
$\phi_j\in\PSH_{\sup}(\alpha)$ with $\phi_j(v)\leq-2^j$.
The finite convex combinations
\[
       s_m=2^{-m}\cdot0+\sum_{j=1}^m2^{-j}\phi_j
\]
are $\alpha$-psh.  This follows first from convexity of the
relative nef cone for model potentials, then by taking decreasing
model approximations in the finitely many slots.
Their coefficients sum to one, $s_m(v_{\triv})=0$, and
$s_{m+1}\leq s_m$ because every $\phi_j\leq0$.
By \Cref{prop:mp-psh-decreasing-stability}, the limit is psh:
its value at $v_{\triv}$ is zero, so it is not identically
$-\infty$.  But $s_m(v)\leq-m$, contradicting non-pluripolarity.
Thus $T(v)<+\infty$.

It follows that $d_\infty(v,w)\leq T(v)+T(w)$.
Symmetry and the triangle inequality hold for each potential
before taking the supremum.  If $d_\infty(v,w)=0$, all
$\alpha$-psh model potentials have the same values at $v$ and $w$.
Every PL function is a scalar multiple of a difference of two
such potentials: use the positive-pair decomposition preceding
\Cref{prop:mp-e1-ma-existence} and divide both potentials by
their common background multiplier.
Since PL functions separate points, $v=w$.
This proves the metric assertion of
\cite[Proposition~3.2.2]{MPWN25}.

\section{Finite-energy measures and dual energy}
\label{sec:mpwn-dual-energy}

For a Radon probability measure $\mu$ on $X^{\bethan}$, define
\[
 E_\alpha^\vee(\mu)
 :=\sup_{\phi\in\mathcal E^1(\alpha)}
       \left\{E_\alpha(\phi)-\int\phi\,\rd \mu\right\}.
\tag{61.6}\label{eq:mpwn-dual-energy}
\]
The expression is invariant under adding constants to $\phi$.  Set
\[
 \mathcal M^1(\alpha)
 :=\{\mu\in\mathcal M_1(X^{\bethan})\mid
                      E_\alpha^\vee(\mu)<+\infty\}.
\tag{61.7}\label{eq:mpwn-finite-energy-measures}
\]
Every $\phi\in\mathcal E^1(\alpha)$ is integrable against every
$\mu\in\mathcal M^1(\alpha)$.  This is what makes every pairing
$\int\phi\,\rd\mu$ below a real number rather than a formal expression,
so it is recorded with its proof.

\begin{lemma}[Finite-energy pairing]
\label{lem:mpwn-integrability}
Let $\phi\in\mathcal E^1(\alpha)$ and $\mu\in\mathcal M^1(\alpha)$.
Then $\int\phi\,\rd\mu$ is finite, and
\[
 E_\alpha(\phi)-E_\alpha^\vee(\mu)\ \leq\ \int\phi\,\rd\mu\ \leq\ \sup\phi .
\]
\end{lemma}

\begin{proof}
A psh function is upper semicontinuous and $X^{\bethan}$ is compact, so
$\phi$ is bounded above and $\int\phi\,\rd\mu$ is well defined in
$[-\infty,+\infty)$, with the stated upper bound since $\mu$ is a
probability measure.  For the lower bound, $\phi$ is an admissible
competitor in the supremum \eqref{eq:mpwn-dual-energy} defining
$E_\alpha^\vee(\mu)$, so
$E_\alpha(\phi)-\int\phi\,\rd\mu\leq E_\alpha^\vee(\mu)$.  Rearranging
gives $\int\phi\,\rd\mu\geq E_\alpha(\phi)-E_\alpha^\vee(\mu)$, and the
right side is finite because $E_\alpha(\phi)>-\infty$ by
\eqref{eq:mp-e1} and $E_\alpha^\vee(\mu)<+\infty$ by
\eqref{eq:mpwn-finite-energy-measures}.
\end{proof}

For $\phi\in\mathcal E^1(\alpha)$, finite energy of
$\MA_\alpha(\phi)$ uses the finite-energy tangent inequality
\cite[Proposition~2.3.5 and its proof]{MPWN25}:
\[
 E_\alpha(\psi)
       \leq E_\alpha(\phi)
               +\int(\psi-\phi)\,\MA_\alpha(\phi),
       \qquad \psi\in\mathcal E^1(\alpha).
\]
The mixed energy integrals here are finite as part of that
finite-energy statement.  Taking the supremum over $\psi$
gives the upper bound below, and the competitor $\psi=\phi$
gives equality:
\[
 E_\alpha^\vee(\MA_\alpha(\phi))
 =E_\alpha(\phi)-\int\phi\,\MA_\alpha(\phi)<+\infty.
\tag{61.8}\label{eq:mpwn-dual-energy-ma}
\]
This establishes membership in $\mathcal M^1(\alpha)$ before
\Cref{lem:mpwn-integrability} is applied to the measure.

\begin{example}
\label{ex:mpwn-divisorial-dirac}
For a divisorial point $v$, monotonicity of energy and
$E_\alpha(c)=c$ give
\[
 E_\alpha^\vee(\delta_v)
 \leq\sup_{\phi\in\mathcal E^1(\alpha)}
          \{\sup\phi-\phi(v)\}<+\infty.
\tag{61.9}\label{eq:mpwn-dirac-energy-bound}
\]
The last bound is the uniform divisorial estimate in
\Cref{thm:mpwn-weak-compactness}.  Thus
$\delta_v\in\mathcal M^1(\alpha)$.
For every $v\in X^{\mathrm{np}}$, the more precise estimate is
\[
              \frac1{n+1}T(v)
              \leq E_\alpha^\vee(\delta_v)\leq T(v).
\tag{61.10}\label{eq:mpwn-point-energy}
\]
We give the truncation argument for
\cite[Proposition~3.2.3]{MPWN25}, whose algebraic antecedent is
\cite[proof of Proposition~11.1, (11.3)]{BJ22}.

First let $\psi\in\PSH(\alpha)$ satisfy
$-M\leq\psi\leq0$ and $\psi(v_{\triv})=0$, with $M\geq0$.
Choose a decreasing model psh net $h_\lambda\downarrow\psi$.
For $0\leq j\leq n$, let $\mu_{\lambda,j}$ be the probability
mixed model measure of $j$ copies of $(\alpha,h_\lambda)$
and $n-j$ copies of $(\alpha,0)$.
Its coefficients are nonnegative mixed intersections of nef
classes on the central components.  Their sum is one, since
intersection with the central fibre can be computed on a
general fibre, where all $n$ classes restrict to $\alpha$.
This is the mixed version of \Cref{prop:mp-model-ma-mass}.
In particular, $\mu_{\lambda,0}=\delta_{v_{\triv}}$.

The model cocycle \eqref{eq:mp-energy-cocycle} and testing formula
\eqref{eq:mp-mixed-testing} give
\[
 E_\alpha(h_\lambda)
     =\frac1{n+1}\sum_{j=0}^n
                      \int h_\lambda\,\rd\mu_{\lambda,j}.
\]
The $j=0$ term is $h_\lambda(v_{\triv})\geq0$ and decreases to
zero.  Each other term is at least $-M$, because
$h_\lambda\geq\psi\geq-M$ and the measures have mass one.
Decreasing continuity of the extended pairing in
\eqref{eq:mp-extended-pairing} therefore yields
\[
                 -\frac{n}{n+1}M
                        \leq E_\alpha(\psi)\leq0.
\tag{61.10a}\label{eq:mpwn-bounded-energy-floor}
\]
The upper bound is monotonicity.  This also proves
$\psi\in\mathcal E^1(\alpha)$.
Only model integrals were used before taking the energy limit;
no unproved convergence of a variable integrand against a
variable measure is required.

If $0<M<T(v)$, choose $\phi\in\PSH_{\sup}(\alpha)$ with
$\phi(v)<-M$, and put $\psi=\max\{\phi,-M\}$.
This maximum is psh by model maximum stability and decreasing
approximation; the real constant can itself be approximated
from above by rational constants.
It satisfies the hypotheses of \eqref{eq:mpwn-bounded-energy-floor}
and $\psi(v)=-M$.  Hence
\[
 E_\alpha^\vee(\delta_v)
       \geq E_\alpha(\psi)-\psi(v)
       \geq\frac{M}{n+1}.
\]
Letting $M\uparrow T(v)$ proves the lower bound in
\eqref{eq:mpwn-point-energy}.  If $T(v)=0$, the same bound follows
by using a constant competitor.
For the upper bound, any finite-energy $\phi$ satisfies
\[
 E_\alpha(\phi)-\phi(v)
                \leq\sup\phi-\phi(v)\leq T(v).
\]
Taking the supremum proves the assertion.  In dimension zero,
connectedness makes $X^{\bethan}$ a single point, and
$T=E_\alpha^\vee(\delta_{v_{\triv}})=0$.
\end{example}

The potential--measure energy gap is
\[
 J_\alpha(\mu,\phi)
 :=E_\alpha^\vee(\mu)-E_\alpha(\phi)+\int\phi\,\rd \mu
 \geq0.
\tag{61.11}\label{eq:mpwn-energy-gap}
\]
It vanishes exactly when $\phi$ maximizes the variational problem for
$\mu$.  For $\mu=\MA_\alpha(\psi)$, it agrees with the asymmetric
Dirichlet quantity $J_\alpha(\psi,\phi)$.

For clarity, the potential-side quantity has the order convention
\[
 J_\alpha(\psi,\phi)
   :=E_\alpha(\psi)-E_\alpha(\phi)
                     +\int(\phi-\psi)\,\MA_\alpha(\psi).
\tag{61.11a}\label{eq:mpwn-potential-energy-gap}
\]
Thus $J_\alpha(\phi)=J_\alpha(0,\phi)$ and
$J_\alpha(\MA_\alpha(\psi),\phi)=J_\alpha(\psi,\phi)$ by
\eqref{eq:mpwn-dual-energy-ma}.  We will use the estimate
\[
 J_\alpha(\phi,\psi)
       \leq C_n\bigl(J_\alpha(\mu,\phi)+J_\alpha(\mu,\psi)\bigr),
\tag{61.11b}\label{eq:mpwn-gap-control}
\]
valid for $\phi,\psi\in\mathcal E^1(\alpha)$ and
$\mu\in\mathcal M^1(\alpha)$, with a constant depending only on
$n$.  This is the mixed-energy input
\cite[Lemma~2.3.6(2)]{MPWN25}.  Its proof first applies the
model estimate, then takes decreasing approximations using
finite-energy pairings.


This Legendre structure is the point at which the model construction becomes
variational.  In the synthetic language of \cite{BJ25FiniteEnergy}, the
energy gap controls both the potential and the measure variables.  What is
specific to the present setting is that the model pairing and the analytic
envelope theorem constructed earlier satisfy the hypotheses needed for that
formalism.

\section{Continuous envelopes}
\label{sec:mpwn-envelopes}

For $f\in C^0(X^{\bethan})$, define
\[
 P_\alpha(f)
 :=\sup\{\phi\in\PSH(\alpha)\mid\phi\leq f\}.
\tag{61.12}\label{eq:mpwn-envelope}
\]
At this point a supremum need be neither upper semicontinuous nor psh, so
the following theorem is substantive.

\begin{theorem}\label{thm:mpwn-envelope-continuity}
For every $f\in C^0(X^{\bethan})$, the function $P_\alpha(f)$ is
continuous and belongs to $\PSH(\alpha)$.
\end{theorem}

\begin{proof}[Imported radial-envelope core]
First, if $\psi\leq f$ is psh and
$\psi_\lambda\downarrow\psi$ is a net in $\mathcal H(\alpha)$,
compactness of $X^{\bethan}$ and continuity of $f$ imply that for every
$\varepsilon\in\mathbf Q_{>0}$, some
$\psi_\lambda-\varepsilon\leq f$. Rational translation keeps this
competitor in $\mathcal H(\alpha)$, so letting $\varepsilon\downarrow0$
gives
\[
 P_\alpha(f)
 =\sup\{\psi\in\mathcal H(\alpha)\mid\psi\leq f\},
\tag{61.13}\label{eq:mpwn-envelope-model-sup}
\]
which is lower semicontinuous.

To prove psh-ness, normalize $f\leq0$ and take the upper envelope of all
complex subgeodesic rays $U\leq0$ whose radial non-Archimedean data obey
$U^{\NA}\leq f$.  Its radial data are psh and bounded above by $f$.
Conversely, each model psh $\psi\leq f$ has a compatible ray and is
therefore dominated by this radial envelope.  Divisorial determination
extends the comparison from $X^{\divis}$ to all of $X^{\bethan}$.
Thus the radial envelope equals $P_\alpha(f)$ and is psh, hence upper
semicontinuous.  Together with lower semicontinuity it is continuous.
The imported radial-envelope assertions are Propositions~2.2.1--2.2.4 of
\cite{MPWN25}; together they identify the two envelopes.  This is the proof
of \cite[Theorem~1.4.1]{MPWN25}.
\end{proof}

\begin{proposition}\label{prop:mpwn-envelope-calculus}
For $f,g\in C^0(X^{\bethan})$ and $c\in\mathbf R$,
\[
 \begin{split}
 P_\alpha(f+c)&=P_\alpha(f)+c,\\
 f\leq g&\Longrightarrow P_\alpha(f)\leq P_\alpha(g),\\
 \|P_\alpha(f)-P_\alpha(g)\|_\infty
       &\leq\|f-g\|_\infty.
 \end{split}
\tag{61.14}\label{eq:mpwn-envelope-calculus}
\]
\end{proposition}

\begin{proof}
The first two statements follow directly from the candidate sets in
\eqref{eq:mpwn-envelope}.  If $\|f-g\|_\infty\leq c$, then
$f\leq g+c$ and $g\leq f+c$; apply the first two statements in both
directions.
\end{proof}

\section{Orthogonality}
\label{sec:mpwn-orthogonality}

\noindent\begin{minipage}{\linewidth}
\begin{theorem}\label{thm:mpwn-orthogonality}
For every $f\in C^0(X^{\bethan})$,
\[
 \int_{X^{\bethan}}
        \bigl(f-P_\alpha(f)\bigr)\,
        \MA_\alpha(P_\alpha(f))=0.
\tag{61.15}\label{eq:mpwn-orthogonality}
\]
Equivalently, the envelope measure is supported on the contact set
$\{P_\alpha(f)=f\}$.
\end{theorem}
\end{minipage}\par

\begin{proof}[Imported restricted-volume core]
First let $f=f_D$ be PL.  Translation of the obstacle translates
its envelope by the same constant and leaves its measure unchanged.
We may therefore assume that $D\geq\cX_0$ on an snc model $\cX$.
The class $A+D$ is big, and
\cite[Proposition~4.2.2]{MPWN25} gives
\[
 \MA_\alpha(P_\alpha(f_D))
  =V^{-1}\sum_{E_i\subset\cX_0}
       b_i\left\langle(A+D)^n\right\rangle_{\cX\mid E_i}
       \delta_{v_{E_i}}.
\tag{61.16}\label{eq:mpwn-envelope-ma-restricted-volume}
\]
The unnormalized masses satisfy
\[
       \sum_i b_i\left\langle(A+D)^n\right\rangle_{\cX\mid E_i}=V
\]
by \cite[Proposition~4.2.1]{MPWN25}; after division by $V$,
the coefficients in \eqref{eq:mpwn-envelope-ma-restricted-volume}
sum to one.

Let $\nu_i=\nu_{E_i}(A+D)$ be the generic minimal multiplicity.
If $\nu_i>0$, the restricted volume along $E_i$ is zero.
For each $i$, the regularization step in \cite[Proposition~4.2.3]{MPWN25}
supplies, for every $\varepsilon>0$, a model psh competitor
$\psi_\varepsilon\leq f_D$ such that
\[
 \psi_\varepsilon(v_{E_i})
        \geq f_D(v_{E_i})-\frac{\nu_i+\varepsilon}{b_i}.
\]
Taking the envelope and then letting $\varepsilon\downarrow0$
gives
\[
 0\leq f_D(v_{E_i})-P_\alpha(f_D)(v_{E_i})
                                      \leq\frac{\nu_i}{b_i}.
\]
In particular, every component with nonzero restricted-volume
mass is a contact point.  This proves orthogonality for a PL
obstacle.  The restricted-volume identity, its mass formula,
the vanishing assertion and the regularization step are the
imported geometric inputs.

For continuous $f$, choose PL functions $f_j\to f$ uniformly.
Put $u_j=P_\alpha(f_j)$ and $u=P_\alpha(f)$.
By \Cref{prop:mpwn-envelope-calculus}, $u_j\to u$ uniformly.
Continuity under uniform convergence
\cite[Proposition~2.3.2]{MPWN25} gives
$\mu_j:=\MA_\alpha(u_j)\rightharpoonup\mu:=\MA_\alpha(u)$.
Set $q_j=f_j-u_j$ and $q=f-u$; these are continuous and
$q_j\to q$ uniformly.  Since the measures have mass one,
\[
 \left|\int q_j\,\rd\mu_j-\int q\,\rd\mu\right|
 \leq\|q_j-q\|_\infty
                   +\left|\int q\,(\rd\mu_j-\rd\mu)\right|
 \longrightarrow0.
\]
The left integrals for $q_j$ vanish by the PL case, proving
\eqref{eq:mpwn-orthogonality}.
This is the deduction of \cite[Theorem~1.5.1]{MPWN25}
from its restricted-volume inputs.
\end{proof}

\begin{remark}\label{rem:mpwn-envelope-first-variation}
For $f,g\in C^0(X^{\bethan})$, orthogonality gives
\[
 \left.\frac{\rd}{\rd t}\right|_{t=0}
 E_\alpha(P_\alpha(f+tg))
                  =\int g\,\MA_\alpha(P_\alpha(f)).
\tag{61.17}\label{eq:mpwn-envelope-differentiability}
\]
Indeed, put $u_t=P_\alpha(f+tg)$ and
$\mu_t=\MA_\alpha(u_t)$.  Each $u_t$ is bounded and psh,
so has finite energy.  The tangent inequality in both orders
gives
\[
 \int(u_t-u_0)\,\rd\mu_t
       \leq E_\alpha(u_t)-E_\alpha(u_0)
       \leq\int(u_t-u_0)\,\rd\mu_0.
\]
On the contact set carrying $\mu_t$, one has
$u_t=f+tg$ and $u_0\leq f$; on the contact set carrying
$\mu_0$, one has $u_0=f$ and $u_t\leq f+tg$.
Hence, for every real $t$,
\[
 t\int g\,\rd\mu_t
       \leq E_\alpha(u_t)-E_\alpha(u_0)
       \leq t\int g\,\rd\mu_0.
\]
The envelope contraction gives
$\|u_t-u_0\|_\infty\leq |t|\|g\|_\infty$, and uniform
convergence of potentials gives $\mu_t\rightharpoonup\mu_0$.
Divide by $t$, reversing the inequalities when $t<0$.
Both one-sided limits equal $\int g\,\rd\mu_0$, which proves
\eqref{eq:mpwn-envelope-differentiability}.
\end{remark}


\section{Strong topologies and quasi-metrics}
\label{sec:mpwn-strong-topologies}

The strong topology on $\mathcal E^1(\alpha)$ is the coarsest topology
finer than weak convergence for which $E_\alpha$ is continuous.  Thus a
net $\phi_i$ converges strongly to $\phi$ exactly when
\[
 \phi_i\to\phi\text{ on }X^{\divis},
 \qquad E_\alpha(\phi_i)\to E_\alpha(\phi).
\tag{61.18}\label{eq:mpwn-potential-strong}
\]
The strong topology on $\mathcal M^1(\alpha)$ is the coarsest topology
finer than weak convergence of measures for which $E_\alpha^\vee$ is
continuous:
\[
 \mu_i\rightharpoonup\mu,
 \qquad E_\alpha^\vee(\mu_i)\to E_\alpha^\vee(\mu).
\tag{61.19}\label{eq:mpwn-measure-strong}
\]

\subsection{Quasi-metrics and the Monge--Amp\`ere map}
\label{subsec:mpwn-quasi-metrics}

We use the quasi-metric convention of \cite[Definition~5.1.1, p.~20]{MPWN25}: a continuous function $d:M\times M\to[0,+\infty)$ vanishes exactly on the diagonal and satisfies
\[
 d(x,z)\leq C\bigl(d(x,y)+d(y,z)\bigr),
 \qquad d(y,x)\leq Dd(x,y)
\]
for all $x,y,z\in M$ and fixed $C,D\geq1$. Its balls induce the specified topology. In particular, symmetry is replaced by comparability of the two orders.

\begin{definition}
\label{def:mpwn-potential-quasidistance}
For $\phi,\psi\in\mathcal E^1(\alpha)$, set
\[
 \partial_\alpha(\phi,\psi)
 :=J_\alpha(\phi,\psi)+|\sup\phi-\sup\psi|.
\tag{61.20}\label{eq:mpwn-partial-distance}
\]
\end{definition}

\begin{theorem}\label{thm:mpwn-strong-quasimetrics}
\begin{enumerate}
\item The function $\partial_\alpha$ is a continuous quasi-metric
      inducing the strong topology on $\mathcal E^1(\alpha)$.
\item There is a unique continuous quasi-metric
      $\delta_\alpha$ on $\mathcal M^1(\alpha)$ inducing the strong
      topology and satisfying
      \[
       \delta_\alpha(\mu,\MA_\alpha(\phi))
               =J_\alpha(\mu,\phi).
      \tag{61.21}\label{eq:mpwn-measure-quasimetric}
      \]
\item The Monge--Amp\`ere map is continuous for the strong topologies.
\end{enumerate}
\par
\end{theorem}

\begin{proof}[Imported quasi-metric theorems]
The first two assertions, including existence of the
measure-side quasi-metric and identification of both topologies,
are \cite[Theorems~5.2.2 and~5.1.2, pp.~20--22]{MPWN25}.
The quasi-metrics may be asymmetric; their quasi-symmetry and
quasi-triangle constants depend only on the dimension.
We retain these results as source inputs.

The quantitative estimate used below is
\cite[Theorem~5.1.2(2), p.~21]{MPWN25}.  For $n\geq1$, put $q=2^{-n}$.
For $\phi,\psi\in\mathcal E^1(\alpha)$,
$\mu,\nu\in\mathcal M^1(\alpha)$, and
\[
 R\geq\max\{1,J_\alpha(\phi),J_\alpha(\psi),
                        E_\alpha^\vee(\mu),E_\alpha^\vee(\nu)\},
\]
it gives
\[
 \left|\int(\phi-\psi)(\rd\mu-\rd\nu)\right|
 \leq C_n J_\alpha(\phi,\psi)^q
                  \delta_\alpha(\mu,\nu)^{1/2}R^{1/2-q}.
\tag{61.21a}\label{eq:mpwn-mixed-gap-estimate}
\]
For example, if $\partial_\alpha(\phi,\psi)=0$, their root
values coincide.  Applying this estimate to $\delta_v$ and
$\delta_{v_{\triv}}$, which have finite energy, gives
$\phi(v)=\psi(v)$ for every divisorial $v$.  Divisorial
determination gives equality everywhere.

If $\phi_i$ converges strongly to $\phi$, continuity of
$\partial_\alpha$ gives $J_\alpha(\phi_i,\phi)\to0$.  Hence
\[
 \delta_\alpha(\MA_\alpha(\phi_i),\MA_\alpha(\phi))
                   =J_\alpha(\phi_i,\phi)\longrightarrow0.
\]
This proves the third assertion from the first two, as in
\cite[Theorem~5.1.3, p.~21]{MPWN25}.  In dimension zero every potential
is constant, $\partial_\alpha$ is absolute difference, and the
measure space has one point; all three assertions hold directly.
\end{proof}

\noindent\begin{minipage}{\linewidth}
\begin{proposition}\label{prop:mpwn-ordered-ma-distances}
Put $\mu_0=\delta_{v_{\triv}}=\MA_\alpha(0)$. For $\mu\in\mathcal M^1(\alpha)$ and $\phi,\psi\in\mathcal E^1(\alpha)$, one has
\[
 \delta_\alpha(\mu,\mu_0)=E_\alpha^\vee(\mu),
 \qquad
 \delta_\alpha(\mu_0,\MA_\alpha(\phi))=J_\alpha(\phi),
\tag{61.21b}\label{eq:mpwn-root-distances}
\]
and
\[
 \begin{aligned}
 \delta_\alpha(\MA_\alpha(\phi),\MA_\alpha(\psi))
     &=J_\alpha(\phi,\psi),\\
 \partial_\alpha(\phi,\psi)
     &=\delta_\alpha(\MA_\alpha(\phi),\MA_\alpha(\psi))
       +|\phi(v_{\triv})-\psi(v_{\triv})|.
 \end{aligned}
\tag{61.21c}\label{eq:mpwn-ma-distance-identity}
\]
Consequently, on the sup-normalized space
\[
 \mathcal E^1_{\sup}(\alpha)
       :=\{\phi\in\mathcal E^1(\alpha):\sup\phi=0\},
\]
the Monge--Amp\`ere map preserves the ordered quasi-distances and is a homeomorphism onto its image for the strong topologies. Moreover, a constant $D_n\geq1$ depending only on $n$ satisfies
\[
 D_n^{-1}J_\alpha(\phi)
       \leq E_\alpha^\vee(\MA_\alpha(\phi))
       \leq D_nJ_\alpha(\phi)
       \qquad(\phi\in\mathcal E^1(\alpha)).
\]
\end{proposition}
\end{minipage}\par

\begin{proof}
Monotonicity gives $E_\alpha(u)\leq\sup u=u(v_{\triv})$, with equality for constants, so $E_\alpha^\vee(\mu_0)=0$. Substituting $0$ or $\mu_0$ in \eqref{eq:mpwn-measure-quasimetric} gives \eqref{eq:mpwn-root-distances}. Substituting $\mu=\MA_\alpha(\phi)$ gives the first equality in \eqref{eq:mpwn-ma-distance-identity} by \eqref{eq:mpwn-dual-energy-ma} and \eqref{eq:mpwn-potential-energy-gap}. The second follows from $\sup\phi=\phi(v_{\triv})$. Quasi-symmetry of $\delta_\alpha$ applied to the two root distances gives the energy comparison.

On $\mathcal E^1_{\sup}(\alpha)$ the root term vanishes. Equality of the measures therefore forces $\partial_\alpha(\phi,\psi)=0$, hence $\phi=\psi$. The same identity identifies convergence of arbitrary nets on the normalized space with convergence of their images to an image point. Surjectivity onto all of $\mathcal M^1(\alpha)$ will be proved in \Cref{thm:mpwn-cy}. For $c\in\mathbf R$ on the unnormalized space, $\MA_\alpha(\phi+c)=\MA_\alpha(\phi)$ while $\partial_\alpha(\phi+c,\phi)=|c|$; this explains the root term.
\end{proof}

\subsection{Joint pairings and point evaluations}
\label{subsec:mpwn-strong-pairings}

\begin{corollary}\label{cor:mpwn-strong-pairing}
For $n\geq1$, put $q=2^{-n}$. Let $\phi,\psi\in\mathcal E^1(\alpha)$ and $\mu,\nu\in\mathcal M^1(\alpha)$, and choose
\[
 R\geq\max\{1,J_\alpha(\phi),J_\alpha(\psi),
                   E_\alpha^\vee(\mu),E_\alpha^\vee(\nu)\}.
\]
Then
\[
 \begin{split}
 \left|\int\phi\,\rd\mu-\int\psi\,\rd\nu\right|
 &\leq |\phi(v_{\triv})-\psi(v_{\triv})|\\
 &\quad+C_n R^{1/2-q}\bigl(
      J_\alpha(\phi,\psi)^q E_\alpha^\vee(\mu)^{1/2}
      +J_\alpha(\psi)^q\delta_\alpha(\mu,\nu)^{1/2}\bigr).
 \end{split}
\tag{61.21d}\label{eq:mpwn-joint-pairing-bound}
\]
Here $C_n$ depends only on $n$. In every dimension the finite-energy pairing $(\phi,\mu)\mapsto\int\phi\,\rd\mu$ is jointly continuous for the strong topologies.
\end{corollary}

\begin{proof}
Every integral is finite by \Cref{lem:mpwn-integrability}. With $\mu_0=\delta_{v_{\triv}}$, write
\[
 \begin{split}
 \int\phi\,\rd\mu-\int\psi\,\rd\nu
 &=\phi(v_{\triv})-\psi(v_{\triv})\\
 &\quad+\int(\phi-\psi)(\rd\mu-\rd\mu_0)
        +\int(\psi-0)(\rd\mu-\rd\nu).
 \end{split}
\]
Apply \eqref{eq:mpwn-mixed-gap-estimate} to the last two terms. For the first, \eqref{eq:mpwn-root-distances} gives $\delta_\alpha(\mu,\mu_0)=E_\alpha^\vee(\mu)$; for the second, $J_\alpha(\psi,0)=E_\alpha^\vee(\MA_\alpha(\psi))\leq D_nJ_\alpha(\psi)$ by \Cref{prop:mpwn-ordered-ma-distances}. Enlarging $C_n$ gives \eqref{eq:mpwn-joint-pairing-bound}.

If $\phi_i\to\phi$ and $\mu_i\to\mu$ strongly, one common $R$ bounds the required energies on a tail of the net: $J_\alpha(\phi_i)=\sup\phi_i-E_\alpha(\phi_i)$ and $E_\alpha^\vee(\mu_i)$ converge. Moreover, the root values converge, $J_\alpha(\phi_i,\phi)\to0$, and $\delta_\alpha(\mu_i,\mu)\to0$. Apply \eqref{eq:mpwn-joint-pairing-bound} with $(\psi,\nu)=(\phi,\mu)$. In dimension zero the potentials are constants and there is only one probability measure, so the conclusion holds directly.
\end{proof}

For $v\in X^{\mathrm{np}}$, recall $T(v)=\sup_{u\in\PSH(\alpha)}(\sup u-u(v))$ from \eqref{eq:mpwn-dinfty}. Take $\phi,\psi\in\mathcal E^1_{\sup}(\alpha)$, $n\geq1$, put $q=2^{-n}$, and choose $R\geq\max\{1,J_\alpha(\phi),J_\alpha(\psi)\}$. Applying \eqref{eq:mpwn-joint-pairing-bound} with $\mu=\nu=\delta_v$ and using \eqref{eq:mpwn-point-energy} gives
\[
 \sup_{\substack{v\in X^{\mathrm{np}}\\T(v)\leq R}}
       |\phi(v)-\psi(v)|
       \leq C_n R^{1-q}J_\alpha(\phi,\psi)^q.
\tag{61.21e}\label{eq:mpwn-bounded-point-energy-evaluation}
\]
The bound $E_\alpha^\vee(\delta_v)\leq T(v)$ makes all evaluations in this display finite. Thus strong convergence of normalized potentials is uniform on each set with bounded $T$.

\begin{example}\label{ex:mpwn-asymmetric-surface-distance}
Let $X=\mathbf P^2$ and $\alpha=c_1(\cO(1))$, so $V=1$. Blow up $(p,0)$ in $X\times\mathbf P^1$, write $E_1$ for the exceptional divisor and $A$ for the pullback of $\alpha$, and choose a rational $0<a<1$ sufficiently small that $A-a[E_1]$ is relatively K\"ahler, as in \Cref{ex:mp-dnc}. Both central multiplicities are one. The model potential $\phi=\phi_{-aE_1}$ is sup-normalized, with values $0$ and $-a$ at the two central vertices $v_{\triv}$ and $v_{E_1}$. The component degrees give
\[
 \mu:=\MA_\alpha(\phi)
       =(1-a^2)\delta_{v_{\triv}}+a^2\delta_{v_{E_1}}.
\]
Since $A|_{E_1}=0$, $A^3=0$, and $E_1^3=1$, the model energy and its measure integral are
\[
 E_\alpha(\phi)=\frac{(A-aE_1)^3}{3}=-\frac{a^3}{3},
 \qquad \int\phi\,\rd\mu=-a^3.
\]
Therefore \eqref{eq:mpwn-root-distances} gives
\[
 \delta_\alpha(\delta_{v_{\triv}},\mu)=J_\alpha(\phi)=\frac{a^3}{3},
 \qquad
 \delta_\alpha(\mu,\delta_{v_{\triv}})
       =E_\alpha^\vee(\mu)=\frac{2a^3}{3}.
\]
Both measures have finite energy, but the two ordered distances differ.
\end{example}


\subsection{Fixed-measure limits and weak convergence}
\label{subsec:mpwn-weak-limits}

\begin{lemma}\label{lem:mpwn-fixed-measure-integral}
Let $\phi_i\in\mathcal E^1(\alpha)$ converge weakly to
$\phi\in\mathcal E^1(\alpha)$, and suppose that $J_\alpha(\phi_i)$
is bounded.  For every $\psi\in\mathcal E^1(\alpha)$ and
$\mu=\MA_\alpha(\psi)$,
\[
                    \int|\phi_i-\phi|\,\rd\mu\longrightarrow0.
\]
\end{lemma}

\begin{proof}
The root values $\sup\phi_i$ converge to $\sup\phi$, so
subtracting them reduces to sup-normalized potentials.
We may also normalize $\psi$.  In dimension zero the assertion
is convergence of these constants, so assume $n\geq1$.
Choose a decreasing model net $\psi_\lambda\downarrow\psi$
and put $\mu_\lambda=\MA_\alpha(\psi_\lambda)$.
Decreasing energy continuity gives strong convergence of
$\psi_\lambda$ to $\psi$; \Cref{thm:mpwn-strong-quasimetrics}
then gives strong convergence of $\mu_\lambda$ to $\mu$.
In particular, their dual energies are eventually bounded.

Since $E_\alpha^\vee(\delta_{v_{\triv}})=0$, applying
\eqref{eq:mpwn-gap-control} with that measure gives
\[
 J_\alpha(\phi_i,\phi)
             \leq C_n\bigl(J_\alpha(\phi_i)+J_\alpha(\phi)\bigr).
\]
The right side is uniformly bounded in $i$.  A single $R$ in
\eqref{eq:mpwn-mixed-gap-estimate} therefore works for all $i$
and all sufficiently large $\lambda$.  Thus
\[
 \left|\int(\phi_i-\phi)\,\rd\mu\right|
 \leq C_R\delta_\alpha(\mu,\mu_\lambda)^{1/2}
                  +\left|\int(\phi_i-\phi)\,\rd\mu_\lambda\right|.
\]
For fixed $\lambda$, the last integral tends to zero because
$\mu_\lambda$ has finite divisorial support and the convergence
of potentials is divisorial.  Choose $\lambda$ first to make
the first term small, then choose $i$ for the second term.
This proves convergence of the signed integrals for nets.

Set $w_i=\max\{\phi_i,\phi\}$.  These are sup-normalized psh
potentials, and $E_\alpha(w_i)\geq E_\alpha(\phi)>-\infty$.
They converge weakly to $\phi$ and have bounded $J_\alpha$.
Apply the signed-integral conclusion also to $w_i$ and use
\[
 |\phi_i-\phi|=2(w_i-\phi)-(\phi_i-\phi)
 \qquad\mu\text{-almost everywhere}.
\]
All terms are integrable by \Cref{lem:mpwn-integrability};
the almost-everywhere qualification excludes their possible
infinite values on null sets.  This proves the assertion.
It reconstructs \cite[Lemma~5.0.5]{MPWN25} with the order of the
two limits explicit.
\end{proof}


\noindent\begin{minipage}{\linewidth}
\begin{warning}{Strong is not weak}
\label{warn:mpwn-strong-not-weak}
A weakly convergent, energy-bounded net need not converge strongly.  The
missing assertion is convergence of its energy, not merely equality of the
energies of possible weak cluster points.  Statements below use the
definitions \eqref{eq:mpwn-potential-strong} and
\eqref{eq:mpwn-measure-strong} without weakening them.
\end{warning}
\end{minipage}\par

\begin{example}\label{ex:mpwn-weak-bounded-not-strong}
On $X=\mathbf P^1$ with $\alpha=c_1(\mathcal O(1))$, let $u_j$
have slope $-1/j$ on the ray $t\ord_0$ for $0\leq t\leq j^2$,
be constant beyond $j^2$, and vanish on every other ray.
The rational scaling in \Cref{ex:mp-scaling-breakpoint} gives the
model psh function $\psi_j(t\ord_0)=-\min\{t,j^2\}$, extended by
zero on the other rays. Thus
$u_j=j^{-1}\psi_j+(1-j^{-1})0$ is a sup-normalized model psh
potential by convexity on a common model. Its curvature measure is
\[
 \mu_j:=\MA_\alpha(u_j)
      =\left(1-\frac1j\right)\delta_{v_{\triv}}
                         +\frac1j\delta_{j^2\ord_0}.
\]
At each fixed divisorial point, $u_j$ tends to zero.  However,
$u_j(j^2\ord_0)=-j$, so the model cocycle
\eqref{eq:mp-energy-cocycle}, with zero root value, gives
\[
 E_\alpha(u_j)=\frac12\int u_j\,\rd\mu_j=-\frac12,
 \qquad J_\alpha(u_j)=\frac12.
\]
Thus the potentials converge weakly with bounded energy, but
do not converge strongly to zero.  On the measure side,
$\mu_j$ tends weakly to $\delta_{v_{\triv}}$: for every bounded
continuous $f$, the difference of its integrals is at most
$2\|f\|_\infty/j$ in absolute value.  Yet
\[
 E_\alpha^\vee(\mu_j)
      =E_\alpha(u_j)-\int u_j\,\rd\mu_j
      =\frac12
      \ne E_\alpha^\vee(\delta_{v_{\triv}})=0.
\]
The same example therefore separates the two measure topologies.
\end{example}


\section{The variational Calabi--Yau theorem}
\label{sec:mpwn-calabi-yau}

In the algebraic setting, \cite[Theorem~12.8]{BJ22} makes surjectivity of the
Monge--Amp\`ere map equivalent to the envelope property and completeness of
the finite-energy space.  The analytic-model envelope and orthogonality
theorems above supply these inputs for a possibly transcendental
K\"ahler class.

\subsection{Existence and the strong inverse}
\label{subsec:mpwn-variational-existence}

\begin{theorem}\label{thm:mpwn-cy}
The probability Monge--Amp\`ere operator is a homeomorphism
\[
 \MA_\alpha:
 \mathcal E^1_{\sup}(\alpha)
       \xrightarrow{\ \sim\ }\mathcal M^1(\alpha)
\tag{61.22}\label{eq:mpwn-cy-homeomorphism}
\]
for the strong topologies.
\end{theorem}

\begin{proof}
If two sup-normalized potentials $\phi,\psi$ have equal
Monge--Amp\`ere measures, then
\[
 J_\alpha(\phi,\psi)+J_\alpha(\psi,\phi)
   =\int(\phi-\psi)
       \bigl(\MA_\alpha(\psi)-\MA_\alpha(\phi)\bigr)=0.
\]
Both summands are nonnegative.  Hence
$\partial_\alpha(\phi,\psi)=0$, and the quasi-metric theorem
gives $\phi=\psi$.

For surjectivity, fix $\mu\in\mathcal M^1(\alpha)$ and choose
a sup-normalized maximizing sequence $u_m\in\mathcal E^1(\alpha)$:
\[
 g_m:=J_\alpha(\mu,u_m)
      =E_\alpha^\vee(\mu)-E_\alpha(u_m)
                            +\int u_m\,\rd\mu\longrightarrow0.
\]
Such a sequence exists because the supremum
\eqref{eq:mpwn-dual-energy} is a finite real number.
Applying \eqref{eq:mpwn-gap-control} to $0,u_m$ gives the
explicit energy bound
\[
 -E_\alpha(u_m)=J_\alpha(u_m)
       \leq C_n\bigl(E_\alpha^\vee(\mu)+g_m\bigr).
\]
After discarding finitely many terms, these energies lie in
one bounded interval.  By \Cref{lem:mpwn-energy-weak-upper},
a subnet $\phi_i=u_{m(i)}$ converges weakly to
$\phi\in\mathcal E^1_{\sup}(\alpha)$.
Compactness is used through a subnet; no sequential compactness
of the weak topology is required.

Set $\mu_m=\MA_\alpha(u_m)$.  Before using surjectivity,
\eqref{eq:mpwn-measure-quasimetric} gives
\[
                   \delta_\alpha(\mu,\mu_m)=g_m\longrightarrow0.
\]
Thus $\mu_m$ converges strongly to $\mu$, and
$E_\alpha^\vee(\mu_m)$ is eventually bounded.
The already established bounds on $J_\alpha(\phi_i)$ and
$J_\alpha(\phi)$, together with \eqref{eq:mpwn-gap-control}
applied to $\delta_{v_{\triv}}$, also bound
$J_\alpha(\phi_i,\phi)$.
For $n\geq1$, \eqref{eq:mpwn-mixed-gap-estimate} consequently gives
\[
 \left|\int(\phi_i-\phi)\,\rd\mu\right|
 \leq C_R g_m^{1/2}
                  +\left|\int(\phi_i-\phi)\,\rd\mu_m\right|,
\]
with $C_R$ independent of $i$ and of sufficiently large $m$.
For fixed $m$, the last term tends to zero by
\Cref{lem:mpwn-fixed-measure-integral}, since
$\mu_m=\MA_\alpha(u_m)$.  Choosing $m$ first and then $i$ proves
\[
                    \int\phi_i\,\rd\mu\longrightarrow\int\phi\,\rd\mu.
\tag{61.23}\label{eq:mpwn-maximizer-integral}
\]
In dimension zero the normalized space has one potential, and
surjectivity holds directly.

Weak upper semicontinuity of energy now yields
\[
 E_\alpha^\vee(\mu)
   \geq E_\alpha(\phi)-\int\phi\,\rd\mu
   \geq\limsup_i\left(E_\alpha(\phi_i)-\int\phi_i\,\rd\mu\right)
   =E_\alpha^\vee(\mu).
\]
Hence $J_\alpha(\mu,\phi)=0$, so
$\MA_\alpha(\phi)=\mu$ by
\eqref{eq:mpwn-measure-quasimetric}.
This gives the variational deduction of
\cite[Propositions~6.0.3--6.0.4]{MPWN25} from the displayed
mixed estimates and the integral lemma.

On sup-normalized potentials,
\[
 \delta_\alpha(\MA_\alpha(\phi),\MA_\alpha(\psi))
       =J_\alpha(\phi,\psi)=\partial_\alpha(\phi,\psi).
\]
The bijection and its inverse are therefore continuous for
the strong topologies.  This is Theorem~B of \cite{MPWN25},
with algebraic predecessor \cite[Theorem~12.8]{BJ22}.
\end{proof}


\subsection{Explicit curve solutions}
\label{subsec:mpwn-curve-solutions}

\begin{example}[The Calabi--Yau solution on $\mathbf P^1$, in closed form]
\label{ex:mpwn-p1-cy}
\Cref{thm:mpwn-cy} is a homeomorphism theorem, and in the smallest
nontrivial case both the solution and every functional attached to it can
be written down.

Let $X=\mathbf P^1$ with $\alpha=c_1(\cO(1))$, so $n=1$ and
$V_\alpha=1$.  Then $X^{\bethan}$ is the trivially valued Berkovich
tree: a root $v_{\triv}$ with one ray $\{t\cdot\ord_x:t\geq0\}$ at each
closed point $x$.  Fix $a\in(0,1)$ and take the divisorial measure
\[
 \mu=a\,\delta_{\ord_0}+(1-a)\,\delta_{\ord_\infty}\in\mathcal M^1(\alpha).
\]

The solution is piecewise linear: let $\phi$ have slope $-a$ along the
ray at $0$ and slope $-(1-a)$ along the ray at $\infty$, each until
$t=1$, and be constant thereafter.  Its Laplacian places
$-(a+(1-a))=-1$ at the root, cancelling
$\MA_\alpha(0)=\delta_{v_{\triv}}$, and places $a$ and $1-a$ at the
two breakpoints, so $\MA_\alpha(\phi)=\mu$.  It is sup-normalized, since
$\phi(v_{\triv})=0$, and
\[
 \phi(\ord_0)=-a,\qquad \phi(\ord_\infty)=-(1-a).
\]

The model cocycle has zero root term,
$\int\phi\,\MA_\alpha(0)=\phi(v_{\triv})=0$. Therefore
\[
 E_\alpha(\phi)=-\tfrac12\bigl(a^2+(1-a)^2\bigr),
 \qquad
 J_\alpha(\phi)=E_\alpha^\vee(\mu)=\tfrac12\bigl(a^2+(1-a)^2\bigr),
\]
the middle equality because $\sup\phi=0$ and the last by
\eqref{eq:mpwn-dual-energy-ma}.

The lower bound in \eqref{eq:mpwn-point-energy} is attained in
dimension one.  For $s\in[0,1]$, let $\phi_s$ have slope $-s$
along the ray at $0$ up to $t=1$, be constant beyond it, and
vanish on every other ray.  Then
\[
 \MA_\alpha(\phi_s)
       =(1-s)\delta_{v_{\triv}}+s\,\delta_{\ord_0},
 \qquad
 \phi_s(\ord_0)=-s,
 \qquad E_\alpha(\phi_s)=-\frac{s^2}{2}.
\]
At $s=1$ the measure is exactly $\delta_{\ord_0}$, so
\eqref{eq:mpwn-dual-energy-ma} evaluates the full variational
supremum:
\[
 E_\alpha^\vee(\delta_{\ord_0})
       =E_\alpha(\phi_1)-\phi_1(\ord_0)=\frac12.
\]

To compute $T(\ord_0)$, one must also bound potentials outside
this family.  For a model psh potential $u$, all outgoing slopes
at the root are nonpositive, by the maximum principle.  The
root atom of its positive curvature measure is one plus the
sum of these slopes.  Thus every root slope is at least $-1$.
Positivity away from the root makes the slope nondecreasing
along each ray, and hence
\[
                    u(\ord_0)\geq u(v_{\triv})-1.
\]
Passing to a decreasing model approximation proves the same
inequality for every $\alpha$-psh potential.  Therefore
$T(\ord_0)\leq1$, while $\phi_1$ attains this bound.
Consequently
\[
       T(\ord_0)=1,\qquad
       E_\alpha^\vee(\delta_{\ord_0})
                         =\frac{T(\ord_0)}{n+1}
                         \quad\text{when }n=1.
\]
This proves sharpness of the lower constant in dimension one.

\end{example}

\begin{example}[The trivial measure]
\label{ex:mpwn-trivial-measure}
The zero potential is sup-normalized and its model measure is
$\delta_{v_{\triv}}$.  Hence the inverse Calabi--Yau homeomorphism sends
$\delta_{v_{\triv}}$ to $0$.  Its dual energy is zero, and its reduced
entropy is also zero because $A_X^{\NA}(v_{\triv})=0$.  This is the
measure-side unit test for both the probability normalization and the
discrepancy shift of \Cref{chap:mp-dual-complexes}.
\end{example}

\begin{remark}[Regularity of the solution]
\label{rem:mpwn-cy-regularity-scope}
The solution belongs to $\mathcal E^1(\alpha)$; it need not be continuous
for an arbitrary finite-energy measure.  Continuity is a further regularity
conclusion under the support hypothesis of
\Cref{thm:mpwn-dual-complex-regularity}.
\end{remark}

\section{Dual complexes and approximation of measures}
\label{sec:mpwn-measure-approximation}

Let $p_{\cX}:X^{\bethan}\to\Delta_{\cX}$ be the retraction and
$i_{\cX}$ the monomial embedding.

\subsection{Restriction and continuity on one dual complex}
\label{subsec:mpwn-convex-restriction}

The uniform divisorial estimate used below is proved in
\Cref{prop:mp-fixed-model-oscillation}; its version-specific source
is \cite[Theorem~4.3.3, pp.~40--41]{MPTropicalV2}.

% RETHLAS-VERIFIED: mpwn-dual-complex-boundary-repair; receipt 20260814T181019Z_3beefd710782
\begin{theorem}\label{thm:mpwn-convex-restriction}
For every \(\phi\in\PSH(\alpha)\) and every snc model \(\cX\),
\[
                      \phi\leq\phi\circ i_{\cX}\circ p_{\cX}.
\tag{61.24}\label{eq:mpwn-retraction-majorant}
\]
Moreover, \(\phi\circ i_{\cX}\) is finite, continuous, and convex on each
face of \(\Delta_{\cX}\).
\end{theorem}
\begin{proof}[Convexity, uniform boundedness, and continuity]
For a model potential, \eqref{eq:mpwn-retraction-majorant} is the
monomial Lelong-number inequality; decreasing approximation proves it in
general.  A relatively K\"ahler model potential is, facewise, a maximum of
finitely many affine functions coming from local generators of a flag
ideal, hence convex.  Decreasing limits preserve convexity in the
extended-real sense.

It remains to justify continuity at the boundary; checking only the finite
vertex set would not suffice.  Let \(\sigma\) be a maximal face and choose a
rational point \(w_\sigma\) in its relative interior.  It is divisorial,
hence a vertex of some snc model \(\cX'\) dominating \(\cX\).  What is needed
here is a bound \emph{uniform in \(\phi\)}, which
\Cref{thm:mp-psh-divisorial} does not state --- its assertions are the
pointwise finiteness \(\phi(v)>-\infty\) and the comparison
\eqref{eq:mp-divisorial-comparison}.  The uniform estimate is the imported
\cite[Theorem~4.3.3]{MP24} recalled in that theorem's proof: it bounds
\(|u(w_\sigma)|\) by a constant \(C_\sigma=C(\cX',\alpha,\gamma)\) for every
normalized model \(\alpha\)-psh \(u\).  Applying it to a decreasing net of
normalized model approximants and passing to the limit gives, after the
normalization \(\sup\phi=0\), a constant \(C_\sigma\) independent of
\(\phi\) with
\[
                       -C_\sigma\leq\phi(w_\sigma)\leq0.
\]
If \(\sigma\) is a point, this already gives the required bound.
Suppose that \(\sigma\) has positive dimension, write its vertices as
\(p_{\sigma,0},\ldots,p_{\sigma,k}\), and write
\[
 w_\sigma=\sum_{i=0}^k a_{\sigma,i}p_{\sigma,i},
 \qquad
 a_\sigma:=\min_i a_{\sigma,i}>0.
\]
For \(w\in\sigma\setminus\{w_\sigma\}\), write
\(w=\sum_i b_i p_{\sigma,i}\), and set
\[
             \lambda:=\min_{b_i>0}\frac{a_{\sigma,i}}{b_i}.
\]
Since \(b_i\leq1\), one has \(\lambda\geq a_\sigma\).  The equality of
the sums of the two systems of barycentric coordinates gives
\(\lambda\leq1\); equality would force \(a_{\sigma,i}=b_i\) for every
\(i\), contrary to \(w\ne w_\sigma\).  Thus
\(a_\sigma\leq\lambda<1\), and
\[
 y:=\frac{w_\sigma-\lambda w}{1-\lambda}
   =\sum_i\frac{a_{\sigma,i}-\lambda b_i}{1-\lambda}
          p_{\sigma,i}
\]
belongs to \(\partial\sigma\): its barycentric coordinates are
nonnegative, sum to one, and one of them vanishes.  In particular,
\(w_\sigma=\lambda w+(1-\lambda)y\).

Convexity in the extended-real sense gives
\[
 -C_\sigma\leq\phi(w_\sigma)
       \leq\lambda\phi(w)+(1-\lambda)\phi(y)
       \leq\lambda\phi(w).
\]
Indeed, the finite left side first rules out \(\phi(w)=-\infty\) and
\(\phi(y)=-\infty\); the last inequality then uses \(\phi(y)\leq0\).
Consequently
\[
             -\frac{C_\sigma}{a_\sigma}\leq\phi(w)\leq0
             \qquad(w\in\sigma).
\]
The same estimate holds at \(w_\sigma\), since \(a_\sigma\leq1\).  For a
zero-dimensional maximal face use the convention \(a_\sigma=1\).

There are only finitely many maximal faces.  Taking the largest of the
facewise constants therefore gives \(M\geq0\) such that every normalized
\(\phi\in\PSH(\alpha)\) satisfies
\[
             -M\leq\phi(w)\leq0
             \qquad(w\in\Delta_\cX).
\]
If \(\psi\in\PSH(\alpha)\) is arbitrary, set
\(\phi=\psi-\sup\psi\).  The preceding estimate shows first that
\(\psi(w)>-\infty\) for every \(w\in\Delta_\cX\); hence
\(\Delta_\cX\subset X^{\mathrm{np}}\).  It also gives, for
\(v,w\in\Delta_\cX\),
\[
             |\psi(v)-\psi(w)|=|\phi(v)-\phi(w)|\leq M.
\]
Taking the supremum over all \(\psi\in\PSH(\alpha)\) proves
\(d_\infty(v,w)\leq M\).  Thus the entire dual complex, not merely its
vertices, is a \(d_\infty\)-bounded subset of \(X^{\mathrm{np}}\).
Proposition~7.1.4 of \cite{MPWN25} applies to this set and proves continuity
on all of \(\Delta_\cX\), including every boundary point.

This supplies the missing boundary step in the printed proof of
\cite[Theorem~7.1.2]{MPWN25}.  The repaired step was independently checked
with verdict \texttt{correct} and no gaps under receipt
\texttt{20260814T181019Z\_3beefd710782}; its exact provenance is recorded in
\Cref{app:provenance}.  The retraction inequality is
\cite[Lemma~7.1.1]{MPWN25}.
\end{proof}

\subsection{Retraction, strong convergence and entropy}
\label{subsec:mpwn-retracted-measures}

For $\mu\in\mathcal M^1(\alpha)$, set
\[
             \mu_{\cX}:=(i_{\cX})_*(p_{\cX})_*\mu.
\tag{61.25}\label{eq:mpwn-retracted-measure}
\]
It is supported on $\Delta_{\cX}$.  As models become finer,
$\mu_{\cX}\rightharpoonup\mu$.

\begin{theorem}\label{thm:mpwn-measure-approximation}
The net $(\mu_{\cX})$ converges strongly to $\mu$.  More precisely,
$E_\alpha^\vee(\mu_{\cX})$ is eventually increasing and
\[
        \lim_{\cX}E_\alpha^\vee(\mu_{\cX})
              =E_\alpha^\vee(\mu).
\tag{61.26}\label{eq:mpwn-retracted-energy}
\]
One can choose a sequence of models $\cX_j$ such that
$\mu_{\cX_j}\to\mu$ strongly and
\[
              \Ent(\mu_{\cX_j})\longrightarrow\Ent(\mu),
\tag{61.27}\label{eq:mpwn-entropy-approximation}
\]
with both sides allowed to be $+\infty$.
\end{theorem}

\begin{proof}
Write $r_{\cX}=i_{\cX}\circ p_{\cX}$.  Every PL function $f$
is represented on some snc model $\cY$, so
$f\circ r_{\cX}=f$ on every model $\cX$ dominating $\cY$.
Consequently
\[
                      \int f\,\rd\mu_{\cX}=\int f\,\rd\mu
\]
eventually.  Uniform density of PL functions in $C^0(X^{\bethan})$,
and the fact that all these measures have mass one, prove
$\mu_{\cX}\rightharpoonup\mu$.

For $\phi\in\PSH(\alpha)$, put $h_{\cX}=\phi\circ r_{\cX}$.
It is finite and continuous by \Cref{thm:mpwn-convex-restriction},
and $\phi\leq h_{\cX}\leq\sup\phi$.
Retraction compatibility in \Cref{lem:mp-retraction} gives
$r_{\cX}\circ r_{\cY}=r_{\cX}$ when $\cY$ dominates $\cX$.
Applying \eqref{eq:mpwn-retraction-majorant} at $r_{\cY}(v)$
then yields $h_{\cY}\leq h_{\cX}$.
The limit is $\phi$: if a model psh function $u\geq\phi$ is
represented on $\mathcal Z$, then $h_{\cX}\leq u\circ r_{\cX}=u$
for every $\cX$ dominating $\mathcal Z$.  Taking decreasing model
approximations to $\phi$ gives $\inf_{\cX}h_{\cX}=\phi$.

The hypotheses of \Cref{cor:usc-net-monotone-integral} now hold
for this decreasing net and the fixed Radon measure $\mu$.
For $\phi\in\mathcal E^1(\alpha)$, it follows that
\[
 \int\phi\,\rd\mu_{\cX}
       =\int h_{\cX}\,\rd\mu
             \downarrow\int\phi\,\rd\mu\in\mathbf R.
\]
In particular,
\[
 E_\alpha(\phi)-\int\phi\,\rd\mu_{\cX}
       \leq E_\alpha(\phi)-\int\phi\,\rd\mu
       \leq E_\alpha^\vee(\mu).
\]
Taking the supremum over $\phi$ first proves that every
$\mu_{\cX}$ has finite energy, bounded by $E_\alpha^\vee(\mu)$.
The same variational expressions increase with the model.
Taking their two suprema in either order gives
\[
 \begin{aligned}
 \sup_{\cX}E_\alpha^\vee(\mu_{\cX})
   &=\sup_{\phi\in\mathcal E^1(\alpha)}
        \sup_{\cX}\left(E_\alpha(\phi)-\int\phi\,\rd\mu_{\cX}\right)\\
   &=E_\alpha^\vee(\mu).
 \end{aligned}
\]
This proves \eqref{eq:mpwn-retracted-energy}.  Together with
weak convergence it proves strong convergence of the model net.

For entropy, the model functions
$a_{\cX}=A_X^{\NA}\circ r_{\cX}$ are nonnegative and continuous,
and increase to $A_X^{\NA}$ by
\eqref{eq:mp-discrepancy-envelope} and
\eqref{eq:mp-reduced-discrepancy}.
Applying \Cref{prop:lsc-net-monotone} therefore gives
\[
 \Ent(\mu_{\cX})=\int a_{\cX}\,\rd\mu
                       \uparrow\Ent(\mu),
\]
including infinite entropy.

To select a sequence, use the quasi-metric rather than only
the numerical energy limits.  Strong convergence of the net
gives $\delta_\alpha(\mu_{\cX},\mu)\to0$.
Choose $\cX_j$ inductively to dominate $\cX_{j-1}$ and to satisfy
\[
 \delta_\alpha(\mu_{\cX_j},\mu)<\frac1j,\qquad
 \begin{cases}
 \Ent(\mu_{\cX_j})>\Ent(\mu)-1/j,
                            &\Ent(\mu)<+\infty,\\
 \Ent(\mu_{\cX_j})>j,        &\Ent(\mu)=+\infty.
 \end{cases}
\]
Each requirement holds on a tail of the directed model system,
so a common dominating model satisfies them simultaneously.
The quasi-metric bound proves strong convergence of this
sequence, while the entropy bounds prove
\eqref{eq:mpwn-entropy-approximation}.
No countable cofinal family of all models is needed.
This reconstructs \cite[Lemmas~7.2.1--7.2.2 and
Corollary~7.2.3]{MPWN25}, with the Radon hypotheses and the
sequential selection explicit.
\end{proof}

\subsection{Finite divisorial interpolation}
\label{subsec:mpwn-atomic-measures}

Use the pushforward formula \eqref{eq:mpwn-retracted-measure} also for arbitrary Radon probability measures. For a fixed snc model $\cX$, entropy of the retracted measure is continuous under weak convergence. Indeed, the function $a_{\cX}=A_X^{\NA}\circ i_{\cX}\circ p_{\cX}$ is continuous and bounded. Thus, for any weakly convergent net of Radon probability measures $\nu_i\rightharpoonup\nu$,
\[
 \Ent((\nu_i)_{\cX})=\int a_{\cX}\,\rd\nu_i
       \longrightarrow\int a_{\cX}\,\rd\nu
       =\Ent(\nu_{\cX}).
\tag{61.27a}\label{eq:mpwn-fixed-model-entropy}
\]
This is the argument of \cite[Lemma~7.2.4, p.~28]{MPWN25}; finiteness of the unretracted measures' energies is not used in this calculation.

Barycentric interpolation on a fixed dual complex supplies finite divisorial support while preserving entropy.

\begin{proposition}\label{prop:mpwn-barycentric-measures}
Fix an snc model $\cX$ and identify its dual complex $\Delta=\Delta_{\cX}$ with its monomial image. Every Radon probability measure $\nu$ supported on $\Delta$ has finite energy and finite entropy. It admits a sequence $\nu_m$ of probability measures with finite divisorial support in $\Delta$ such that
\[
 \begin{gathered}
 \nu_m\longrightarrow\nu\quad\text{strongly},\qquad
 E_\alpha^\vee(\nu_m)\uparrow E_\alpha^\vee(\nu),\\
 \Ent(\nu_m)=\Ent(\nu)\quad\text{for every }m.
 \end{gathered}
\tag{61.27b}\label{eq:mpwn-barycentric-convergence}
\]
\end{proposition}

\begin{proof}
The uniform bound in the proof of \Cref{thm:mpwn-convex-restriction} gives $0\leq M_\Delta<+\infty$ such that $-M_\Delta\leq u|_\Delta\leq0$ for every sup-normalized $u\in\PSH(\alpha)$. In particular $T(v)\leq M_\Delta$ for $v\in\Delta$. For normalized $u\in\mathcal E^1(\alpha)$, monotonicity gives $E_\alpha(u)\leq0$, so
\[
 0\leq E_\alpha^\vee(\nu)
      =\sup_{u\in\mathcal E^1_{\sup}(\alpha)}\left(E_\alpha(u)-\int u\,\rd\nu\right)
      \leq M_\Delta.
\]
The discrepancy $A_X^{\NA}|_\Delta$ is finite, continuous and affine on each face by \eqref{eq:mp-total-discrepancy} and \eqref{eq:mp-reduced-discrepancy}. It is therefore bounded on $\Delta$, proving finite entropy.

Choose nested finite rational triangulations $\mathcal T_m$ of $\Delta$, each subordinate to its original faces, with mesh tending to zero. Successive barycentric subdivisions of one compatible rational triangulation suffice. For each vertex $v$ of $\mathcal T_m$, let $\lambda_{m,v}$ be its continuous piecewise affine nodal function, taking value $1$ at $v$ and $0$ at every other vertex. On every simplex these functions are its barycentric coordinates; hence they are nonnegative and sum to one. Define
\[
 \nu_m:=\sum_{v\in\operatorname{Vert}(\mathcal T_m)}
                  \left(\int_\Delta\lambda_{m,v}\,\rd\nu\right)\delta_v.
\tag{61.27c}\label{eq:mpwn-barycentric-measure}
\]
Every vertex is rational and hence divisorial. The coefficients are nonnegative real numbers of sum one; zero coefficients may be omitted.

For $g\in C^0(\Delta)$, put $I_mg=\sum_v g(v)\lambda_{m,v}$. Then
\[
                   \int g\,\rd\nu_m=\int I_mg\,\rd\nu.
\]
At each point, $I_mg$ is a convex combination of values of $g$ at vertices of its simplex. Uniform continuity of $g$ and the vanishing mesh therefore give $\|I_mg-g\|_\infty\to0$. Applying this to restrictions of continuous functions on $X^{\bethan}$ proves $\nu_m\rightharpoonup\nu$.

For $u\in\mathcal E^1(\alpha)$, its restriction to each face is continuous and convex by \Cref{thm:mpwn-convex-restriction}. Convexity gives $u\leq I_mu$. Since every refined simplex lies in a simplex on which $I_mu$ is affine, the same barycentric comparison gives $I_{m+1}u\leq I_mu$. Thus $I_mu\downarrow u$ uniformly on $\Delta$, and
\[
 E_\alpha(u)-\int u\,\rd\nu_m
       \uparrow E_\alpha(u)-\int u\,\rd\nu.
\]
Taking the two suprema in either order gives
\[
 \sup_m E_\alpha^\vee(\nu_m)
   =\sup_{u\in\mathcal E^1(\alpha)}\sup_m
       \left(E_\alpha(u)-\int I_mu\,\rd\nu\right)
   =E_\alpha^\vee(\nu).
\]
This proves the energy assertion in \eqref{eq:mpwn-barycentric-convergence}. Together with weak convergence this proves strong convergence. Finally $I_m(A_X^{\NA}|_\Delta)=A_X^{\NA}|_\Delta$, because the discrepancy is affine on every original face. Integrating gives the exact entropy identity. The construction also applies to zero-dimensional faces, where interpolation leaves every function unchanged.
\end{proof}

\begin{example}\label{ex:mpwn-interval-barycentric-measures}
Take $X=\mathbf P^1$ and $\alpha=c_1(\cO(1))$. On the model obtained by blowing up $(p,0)$, the dual interval consists of $v_t=t\ord_p$ for $0\leq t\leq1$, with $v_0=v_{\triv}$. The reduced discrepancy is affine on this interval. Its root value is $0$ and its exceptional value is $1$, since the exceptional divisor has central multiplicity $1$ and relative canonical coefficient $1$. Thus $A_X^{\NA}(v_t)=t$, consistently with \Cref{ex:mp-discrepancy-curve-normalization}. Let $\nu$ be the image of Lebesgue probability measure on $[0,1]$. For $N=2^m$ and $h=N^{-1}$, \eqref{eq:mpwn-barycentric-measure} becomes
\[
 \nu_m=\frac h2\delta_{v_0}
              +h\sum_{k=1}^{N-1}\delta_{v_{kh}}
              +\frac h2\delta_{v_1}.
\]
The endpoint hat functions have integral $h/2$ and the interior ones have integral $h$. These weights sum to one, and
\[
 \Ent(\nu_m)=h\sum_{k=1}^{N-1}kh+\frac h2
            =\frac12=\Ent(\nu).
\]
Thus a measure with no atoms is approximated strongly by finite divisorial measures on the same interval, with exactly the same entropy.

Its normalized Calabi--Yau solution can also be computed. On the ray
$t\ord_p$ put
\[
 \begin{gathered}
 u(t\ord_p)=
 \begin{cases}
 -t+t^2/2,&0\leq t\leq1,\\
 -1/2,&t\geq1,
 \end{cases}
 \\
 u=0\quad\text{on every other ray}.
 \end{gathered}
\tag{61.27d}\label{eq:mpwn-interval-continuous-solution}
\]
At the supported endpoint of the ray take the same value $-1/2$.
\par\noindent\begin{minipage}{\linewidth}
For the atomic measure $\nu_m$, define
\[
 \begin{gathered}
 u_m(t\ord_p)=-\int_0^1\min\{t,x\}\,\rd\nu_m(x),\\
 u_m=0\quad\text{on every other ray}.
 \end{gathered}
\]
\end{minipage}\par
For each positive rational $x$, the summand $-\min\{t,x\}$,
extended by zero on the other rays, is the rational scaling of
the model psh function $-\min\{t,1\}$ in
\Cref{ex:mp-scaling-breakpoint}. At $x=0$ it is the zero potential.
Thus $u_m$ is a convex combination of rational model psh functions;
on one common model, nefness of their classes proves the same
property for $u_m$. In dimension one the intersection formula is
affine in the vertical divisor, and gives $\MA_\alpha(u_m)=\nu_m$.
The initial slope is $-(1-h/2)$, leaving root mass $h/2$.
Each interior breakpoint increases the slope by $h$, and the
breakpoint at $1$ increases it by $h/2$ to zero. These slopes
check the same atomic measure directly. The potentials are
sup-normalized and rational PL.

At every grid point $kh$, integration of the piecewise affine
function $x\mapsto\min\{kh,x\}$ by the displayed trapezoidal weights
is exact. Consequently $u_m$ is the affine interpolant of $u$
on this grid. On $[kh,(k+1)h]$ the interpolation error is
\[
 u_m(t\ord_p)-u(t\ord_p)
       =\frac12(t-kh)((k+1)h-t).
\]
The interpolants decrease under dyadic refinement. Decreasing
stability gives the bounded psh limit $u$, and hence
\[
 u_m\downarrow u,\qquad
 \|u_m-u\|_\infty=\frac{h^2}{8},\qquad
 \MA_\alpha(u)=\nu.
\tag{61.27e}\label{eq:mpwn-interval-uniform-error}
\]
The last equality follows from continuity under uniform convergence
of potentials, \cite[Proposition~2.3.2]{MPWN25}.
Thus $u$ is continuous and psh, while its quadratic restriction
to the interval shows that it is not PL.

Finally, $\int_0^1u(t\ord_p)\,\rd t=-1/3$. The trapezoidal error
for this quadratic is $h^2/12$, so the model cocycle gives
\[
 \begin{gathered}
 E_\alpha(u_m)=-\frac16+\frac{h^2}{24},\qquad
 E_\alpha(u)=-\frac16,\\
 E_\alpha^\vee(\nu_m)=J_\alpha(u_m)=\frac16-\frac{h^2}{24},\\
 E_\alpha^\vee(\nu)=J_\alpha(u)=\frac16 .
 \end{gathered}
\tag{61.27f}\label{eq:mpwn-interval-exact-energies}
\]
The limiting formulas follow by decreasing energy continuity and
\eqref{eq:mpwn-dual-energy-ma}; all potentials are bounded, so have
finite energy. This calculation checks the strong limit as well
as the exact entropy value $1/2$.
\end{example}


\section{Comparison, locality, and regularity}
\label{sec:mpwn-regularity}

\subsection{Locality and support domination}
\label{subsec:mpwn-locality}

\begin{theorem}\label{thm:mpwn-locality}
For $\phi,\psi\in\mathcal E^1(\alpha)$ and
$U=\{\phi>\psi\}$,
\[
 \mathbf1_U\MA_\alpha(\max\{\phi,\psi\})
      =\mathbf1_U\MA_\alpha(\phi).
\tag{61.28}\label{eq:mpwn-locality}
\]
\end{theorem}

\begin{proof}[Imported locality theorem]
This is \cite[Theorem~8.1.1, pp.~28--29]{MPWN25}.
Its model core represents $\phi,\psi$ and $\max\{\phi,\psi\}$ by
vertical divisors on one snc model. Along a central component whose
vertex satisfies $\phi>\psi$, the divisors representing $\phi$ and
the maximum agree in a neighborhood of that component. Their
restrictions and atomic intersection masses therefore agree.

The finite-energy conclusion is retained as part of the imported
theorem. Decreasing model approximants give weakly convergent moving
measures, but weak convergence alone does not justify multiplication
by the moving differences $\max\{\phi_j-\psi_j,0\}$. The algebraic
antecedent \cite[Theorem~7.40 and its proof, pp.~768--770]{BJ22}
uses the weighted continuity statement of its Proposition~7.19(iii)
for this step. The cited transcendental proof summarizes the passage
by decreasing continuity; we retain its full locality assertion as
the source input, rather than claim an independent reconstruction
from that summary.
\end{proof}

The following domination principle will also be used in \Cref{chap:mpwn-k-stability}.

\begin{corollary}[Domination on the support]
\label{cor:mpwn-domination}
Let $\phi,\psi\in\mathcal E^1(\alpha)$ and put
$K=\operatorname{supp}\MA_\alpha(\phi)$.  If $\phi\geq\psi$ on $K$,
then $\phi\geq\psi$ on all of $X^{\bethan}$.
\end{corollary}

\begin{proof}
Fix $\varepsilon>0$, put $\chi:=\max\{\phi+\varepsilon,\psi\}$ and
$U:=\{\phi+\varepsilon>\psi\}$.  Constant invariance and
\Cref{thm:mpwn-locality} give
\[
 \mathbf1_U\MA_\alpha(\chi)=\mathbf1_U\MA_\alpha(\phi).
\]
The measure $\mu:=\MA_\alpha(\phi)$ has finite energy.
By \Cref{lem:mpwn-integrability}, both $\phi$ and $\psi$ are finite
$\mu$-almost everywhere.  At every such point of $K$, the hypothesis
$\phi\geq\psi$ implies $\phi+\varepsilon>\psi$.
Thus $\mu(K\setminus U)=0$ and $\mu(U)=1$.
The pointwise inclusion $K\subseteq U$ is unnecessary and need not
hold where both potentials equal $-\infty$.

Locality now implies $\MA_\alpha(\chi)(U)=1$.
Both Monge--Amp\`ere measures are probabilities, so
\[
 \MA_\alpha(\chi)=\mathbf1_U\MA_\alpha(\chi)
                 =\mathbf1_U\mu=\mu.
\]
Injectivity in \Cref{thm:mpwn-cy} gives $\chi=\phi+c$ for a
constant $c$.  On the full $\mu$-measure subset of $K$ where
$\phi$ and $\psi$ are finite, we also have $\chi=\phi+\varepsilon$.
Evaluation at any point of this subset gives $c=\varepsilon$.
Hence $\psi\leq\phi+\varepsilon$ everywhere.
Letting $\varepsilon\downarrow0$ proves the assertion, including
at points where $\phi=-\infty$.
\end{proof}

\subsection{Continuity of solutions}
\label{subsec:mpwn-continuous-solutions}

\begin{theorem}\label{thm:mpwn-dual-complex-regularity}
Let $\mu\in\mathcal M^1(\alpha)$.  If
$\operatorname{supp}\mu\subset\Delta_{\cX}$ for one snc model, then the
sup-normalized solution of
$\MA_\alpha(\phi)=\mu$ is continuous on $X^{\bethan}$.
\end{theorem}

\begin{proof}
The restriction $\phi|_{\Delta_{\cX}}$ is continuous by
\Cref{thm:mpwn-convex-restriction}.  Choose a net
$\phi_i\in\mathcal H(\alpha)$ decreasing pointwise to $\phi$.
On the compact dual complex, the net form of Dini's theorem gives
uniform convergence.  Thus, for every $\varepsilon>0$, all sufficiently
late approximants satisfy
\[
              \phi\leq\phi_i\leq\phi+\varepsilon
       \quad\text{on }\operatorname{supp}\mu.
\]
Apply \Cref{cor:mpwn-domination} with reference potential $\phi$
and comparison potential $\phi_i-\varepsilon$.
It gives $\phi_i-\varepsilon\leq\phi$ everywhere, while
$\phi\leq\phi_i$ already holds by the decreasing approximation.
Consequently $\phi_i\to\phi$ uniformly on $X^{\bethan}$, and
$\phi$ is continuous.  This is \cite[Theorem~8.2.1]{MPWN25}.
\end{proof}

\noindent\begin{minipage}{\linewidth}
\begin{warning}{Continuity is not model regularity}
\label{warn:mpwn-continuous-not-model}
The theorem proves $\phi\in C^0\cap\PSH(\alpha)$.  It does not say that
$\phi$ is PL, that it is determined by one model, or that it is constant
on the fibers of $p_{\cX}$.  Convexity of its restriction to
$\Delta_{\cX}$ is a separate statement.
\end{warning}
\end{minipage}\par

\subsection{Envelopes and entropy-preserving density}
\label{subsec:mpwn-envelope-density}

\begin{corollary}\label{cor:mpwn-divisorial-envelope}
If $\mu=\sum_{i=1}^Na_i\delta_{v_i}$ is a probability measure with
$v_i\in X^{\divis}$, then there exists
$f\in\PL_{\mathbf R}(X^{\bethan})$ such
that
\[
                   \MA_\alpha(P_\alpha(f))=\mu.
\tag{61.29}\label{eq:mpwn-divisorial-envelope-solution}
\]
\end{corollary}

\begin{proof}
Put all $v_i$ among the vertices of one snc model, and write its entire
vertex set as $\{w_j=v_{E_j}\}_{j=1}^m$, with
$\cX_0=\sum_{j=1}^m b_jE_j$.  The solution
$\phi$ is continuous by
\Cref{thm:mpwn-dual-complex-regularity}.  Define the vertical real divisor
\[
                     G:=\sum_{j=1}^m b_j\phi(w_j)E_j
\]
and let $f=\phi_G\in\PL_{\mathbf R}$; thus
$f(w_j)=\phi(w_j)$ at every vertex, not only at the support vertices.
The finite-constraint envelope lemma \cite[Lemma~7.1.5]{MPWN25} gives
\[
 P_\alpha(f)=\sup\{u\in\PSH(\alpha):u(w_j)\leq\phi(w_j)
                                      \text{ for every }j\}.
\]
Hence $\phi\leq P_\alpha(f)$.  On the other hand,
$P_\alpha(f)(v_i)\leq f(v_i)=\phi(v_i)$ at every point of
$\operatorname{supp}\mu$; the domination consequence of locality gives
$P_\alpha(f)\leq\phi$ everywhere.  Thus equality holds, and its
Monge--Amp\`ere measure is $\mu$.  This is
\cite[Proposition~8.3.2]{MPWN25}.
\end{proof}

\noindent\begin{minipage}{\linewidth}
\begin{corollary}\label{cor:mpwn-fixed-complex-uniform-solutions}
Fix an snc model $\cX$ and put $\Delta=\Delta_{\cX}$. For $n\geq1$, there is a constant $C_{\alpha,\Delta}$ such that the sup-normalized solutions $\phi,\psi$ for any two Radon probability measures $\mu,\nu$ supported on $\Delta$ satisfy
\[
 \|\phi-\psi\|_{C^0(X^{\bethan})}
        \leq C_{\alpha,\Delta}\,
                         \delta_\alpha(\mu,\nu)^{2^{-n}}.
\tag{61.29a}\label{eq:mpwn-fixed-complex-uniform-bound}
\]
For every such $\mu$, its solution is a uniform and strong limit of normalized envelopes $P_\alpha(f_m)$ with $f_m\in\PL_{\mathbf R}(X^{\bethan})$, whose Monge--Amp\`ere measures have finite divisorial support in $\Delta$ and entropy equal to $\Ent(\mu)$.
\end{corollary}
\end{minipage}\par

\begin{proof}
All the measures have finite energy by \Cref{prop:mpwn-barycentric-measures}, so the solutions exist by \Cref{thm:mpwn-cy} and are continuous by \Cref{thm:mpwn-dual-complex-regularity}. Use the constant $M_\Delta$ from the proof of \Cref{prop:mpwn-barycentric-measures}. Then $T|_\Delta\leq M_\Delta$ and both dual energies are at most $M_\Delta$. By \Cref{prop:mpwn-ordered-ma-distances}, both potential energies $J_\alpha$ are at most $D_nM_\Delta$. Set $R=\max\{1,D_nM_\Delta\}$ and $q=2^{-n}$. Equations \eqref{eq:mpwn-bounded-point-energy-evaluation} and \eqref{eq:mpwn-ma-distance-identity} give
\[
 \sup_\Delta|\phi-\psi|
       \leq C_nR^{1-q}\delta_\alpha(\mu,\nu)^q=:\varepsilon.
\]
On $\operatorname{supp}\mu$ one has $\phi\geq\psi-\varepsilon$, so \Cref{cor:mpwn-domination} gives this inequality everywhere. Applying the same principle with reference potential $\psi$ gives $\psi\geq\phi-\varepsilon$ everywhere. This proves \eqref{eq:mpwn-fixed-complex-uniform-bound}, with $C_{\alpha,\Delta}=C_nR^{1-q}$ independent of the measures.

Apply \Cref{prop:mpwn-barycentric-measures} to $\mu$. The solutions for its approximants converge strongly by \Cref{thm:mpwn-cy}, and uniformly by \eqref{eq:mpwn-fixed-complex-uniform-bound}. Each is an envelope of a real PL obstacle by \Cref{cor:mpwn-divisorial-envelope}; translating the obstacle normalizes its envelope. Their entropies equal $\Ent(\mu)$ by construction. In dimension zero all normalized solutions are zero, and the approximation statement holds with a constant sequence.
\end{proof}

\begin{corollary}\label{cor:mpwn-divisorial-entropy-density}
For every $\phi\in\mathcal E^1_{\sup}(\alpha)$ there are $f_j\in\PL_{\mathbf R}(X^{\bethan})$ such that $\phi_j=P_\alpha(f_j)$ is sup-normalized, $\MA_\alpha(\phi_j)$ has finite divisorial support, and
\[
 \phi_j\longrightarrow\phi\quad\text{strongly},\qquad
 \Ent(\MA_\alpha(\phi_j))
           \longrightarrow\Ent(\MA_\alpha(\phi)).
\tag{61.29b}\label{eq:mpwn-divisorial-entropy-density}
\]
The entropy limit may be infinite.
\end{corollary}

\begin{proof}
Put $\mu=\MA_\alpha(\phi)$. Choose the retracted measures $\rho_j=\mu_{\cX_j}$ from \Cref{thm:mpwn-measure-approximation}, so $\rho_j\to\mu$ strongly and $\Ent(\rho_j)\to\Ent(\mu)$. On each fixed $\Delta_{\cX_j}$ apply \Cref{prop:mpwn-barycentric-measures} and select a finite divisorial measure $\nu_j$ satisfying
\[
 \delta_\alpha(\nu_j,\rho_j)<1/j,
 \qquad \Ent(\nu_j)=\Ent(\rho_j).
\]
The quasi-triangle inequality gives $\delta_\alpha(\nu_j,\mu)\to0$. The Calabi--Yau homeomorphism and \Cref{cor:mpwn-divisorial-envelope} then give normalized $\phi_j=P_\alpha(f_j)$ converging strongly to $\phi$. The displayed entropy identity proves \eqref{eq:mpwn-divisorial-entropy-density}, including the infinite limit. The models may vary with $j$, so this argument yields strong convergence; the uniform estimate \eqref{eq:mpwn-fixed-complex-uniform-bound} applies when the supports lie in one fixed dual complex.
\end{proof}


\section*{Source note}
\label{sec:mpwn-cy-source-note}

\enlargethispage{2\baselineskip}

Weak compactness and non-pluripolar points are MPWN25 Section~3.
Continuity of envelopes and orthogonality are Section~4.  Strong
quasi-metrics and continuity of Monge--Amp\`ere are Section~5, and the
Calabi--Yau homeomorphism is Section~6.  Its algebraic variational
antecedent is \cite[Theorem~12.8]{BJ22}, while the mixed-energy estimates
also fit the natural synthetic framework of \cite{BJ25FiniteEnergy}.
Dual-complex convexity and
measure approximation are Section~7.  Locality and the regularity theorem
are Section~8; compare \cite[Theorems~7.40, 11.12, and~12.12]{BJ22} for the
algebraic locality, bounded-set continuity, and regularity pattern.
The published synthetic measure construction is
\cite[Theorems~2.22, 2.23 and~2.25, pp.~182--185]{BJ25FiniteEnergy}:
orthogonality gives strong density of model measures before
surjectivity, and the ordered Dirichlet quasi-metric is then
constructed by maximizing sequences. The proofs use the estimates
of Boucksom--Jonsson, while the analytic-model and restricted-volume
inputs retain their stated transcendental scope. Locality is retained
as the full imported assertion of MPWN25 Theorem~8.1.1; its summary
of decreasing continuity is not offered as a new weighted-measure
limit proof.  The subnet correction in weak compactness and the
whole-complex boundedness repair in
\Cref{thm:mpwn-convex-restriction} are displayed explicitly rather than
silently attributed to the printed proofs.
